% !TEX root = ../../main.tex
% agent_8_5d_hcs_cfg.tex
% Chain-level BV derivation of the universal Phi-trace identity from
% 5d holomorphic Chern--Simons on R_t x C^2 (Costello--Yagi 2018),
% read side-by-side with Costello--Francis--Gwilliam 2026. Wave-file
% workspace; reader-facing inscription lives in
% chapters/theory/phi_universal_trace_platonic.tex.

\section*{Agent 8. The 5d holomorphic Chern--Simons realisation of \texorpdfstring{$\PhiFA_3$}{Phi-FA-3}}
\label{sec:wave-cfg-agent-8}

\subsection*{The theorem the identity rests on}

The Universal $\Phi$-Trace Identity of
Theorem~\ref{thm:phi-universal-trace} pairs two supertraces on one
centre $\mathfrak{Z}(\mathcal{A}) = \int_{S^1 \times \Ran(X)}\mathcal{A}$.
Both supertraces are partition-function integrals in a BV field theory;
the field theory that computes them at Stage~$1$ is $5$d holomorphic
Chern--Simons on $Y = \R_t \times \C^2$ with a simply-laced gauge Lie
algebra $\frakg$, and the BV partition function is the Feynman-graph
expansion of the action
\begin{equation}
\label{eq:5d-hCS-action}
  S_{5\mathrm{d\,hCS}}[\cA] \;=\; \int_{Y}\,
  \Omega_{\C^2}\wedge \operatorname{CS}_3\bigl(\cA;\, d_t + \bar\partial_{\C^2}\bigr),
\qquad
  \cA \in \Omega^{0,\bullet}(\C^2)\otimes\Omega^\bullet(\R_t)\otimes \frakg[1],
\end{equation}
with $\Omega_{\C^2} = dz^1\wedge dz^2$ the flat holomorphic top-form.
The cubic interaction is
$\operatorname{CS}_3(\cA) = \tfrac{1}{2}\cA\wedge(\bar\partial_{\C^2}+d_t)\cA
+ \tfrac{1}{6}\cA\wedge[\cA,\cA]$; the classical BV antibracket comes
from the pairing
$(-,-)_{BV} : \Omega^{0,\bullet}(\C^2) \otimes \Omega^\bullet(\R_t)
\otimes\frakg[1]\; \otimes\; \Omega^{2,2-\bullet}(\C^2)
\otimes\Omega^{1-\bullet}(\R_t)\otimes\frakg[1] \;\to\; \C$
of total ghost degree $-1$, namely $\int \mathrm{tr}_\frakg(\alpha\wedge\beta)
\wedge \Omega_{\C^2}$ with $\mathrm{tr}_\frakg$ the Killing form.

CFG 2026 establishes the parallel statement in purely topological
dimension: $3$d ordinary Chern--Simons has, as its $E_3$-factorisation
algebra of local observables, a presentation equivalent to the
factorisation homology of topological CS on the handlebody stratification
(Costello--Francis--Gwilliam 2026~\cite{CFG2026}). Side-by-side, the
present wave file asserts and proves, at the chain level, three
claims which together pin Stage~$1$ of the two-stage factorisation
$\Phi_3 = \SpCh_{\Sigma_2,C}\circ\PhiFA_3$ to the $5$d hCS partition
function.

\subsection*{Three theorems, one BV integral}

\paragraph{Theorem~A (Stage-$1$ identification).}
For $X = \C^3$ presented as the flat holomorphic part $\C^2$ of
$Y = \R_t\times\C^2$ with topological $\R_t$, the Stage-$1$
factorisation algebra $\PhiFA_3(\Perf(\C^3))$ is equivalent, at the
chain level, to the BV observables of $5$d hCS on $Y$ with gauge
$\frakg = \widehat{\mathfrak{gl}}_1$, Wick-rotated to the topological
limit along $\R_t$:
\begin{equation}
\label{eq:stage1-5d-hCS}
  \PhiFA_3\bigl(\Perf(\C^3)\bigr) \;\simeq\;
  \Obs^{\mathrm{BV}}_{5\mathrm{d\,hCS}}\bigl(\C^2\times\R_t;\, \widehat{\mathfrak{gl}}_1\bigr)\Big|_{\hbar\to 0}
  \;\otimes_{\widehat{\mathfrak{gl}}_1}\; \End\bigl(\mathcal{F}_{\mathrm{Fock}}\bigr).
\end{equation}

\paragraph{Theorem~B (universal trace = BV pairing).}
The universal trace of
Theorem~\ref{thm:phi-universal-trace} is the
BV partition function pairing
\begin{equation}
\label{eq:universal-trace-BV}
  \tr_{\mathfrak{Z}(\mathcal{A}_X)}\bigl(\mathfrak{K}_{\mathcal{A}_X}\bigr)
  \;=\;
  \bigl\langle e^{S_{5\mathrm{d\,hCS}}/\hbar}\bigr\rangle_{\mathrm{BV}}^{\mathfrak{K}}
  \;=\;
  \int_{\mathcal{L}_{\mathrm{BV}}} e^{S_{5\mathrm{d\,hCS}}[\cA]/\hbar}\,
  \mathrm{Ch}_{\mathfrak{K}}(\cA)\,D\cA,
\end{equation}
evaluated on the BV Lagrangian $\mathcal{L}_{\mathrm{BV}}\subset
\Fields_{\mathrm{BV}}(Y,\frakg)$ of gauge-fixed Costello--Li heat-kernel
regularisation, paired against the Chern character
$\mathrm{Ch}_{\mathfrak{K}}(\cA)$ of the Koszul involution on the
off-shell field content. On the Borcherds side, the same BV integral
specialised to the framed $K3\times E$ Stage-$2$ curve $C\subset \C^2$
computes $\kappa_{\mathrm{BKM}}(\mathfrak{g}_{\Delta_5}) = c_1(0)/2 = 5$
as the one-loop singular-theta correspondence residue of the BV
partition function on the Heegner-lift locus.

\paragraph{Theorem~C (all-orders convergence for simply-laced).}
For simply-laced $\frakg$, the $\hbar$-expansion of the BV partition
function~\eqref{eq:universal-trace-BV} converges (is not merely asymptotic):
\begin{equation}
\label{eq:all-orders-convergence}
  \bigl\langle e^{S/\hbar}\bigr\rangle_{\mathrm{BV}}
  \;=\; \sum_{n\ge 0}\hbar^n\,
  Z^{(n)}_{5\mathrm{d\,hCS}}(\frakg)
  \quad\text{with radius of convergence}\quad R(\frakg) \;\ge\; \rho_\frakg > 0,
\end{equation}
with $\rho_\frakg$ the inverse Kontsevich-weight entropy of the
Feynman-graph sum on $\C^2$ (finite by Kontsevich formality).

\subsection*{The three-paragraph proof (chain-level BV)}

\paragraph{Step~1 (BV complex as tensor factorisation).}
The classical BV complex is
\begin{equation}
\label{eq:bv-complex-tensor}
  \Obs^{\mathrm{cl}}_{5\mathrm{d\,hCS}}(\R_t\times\C^2;\frakg)
  \;=\; \Bigl( \Omega^\bullet_c(\R_t) \otimes \Omega^{0,\bullet}_c(\C^2)\otimes \frakg[1],\,
  Q_{\mathrm{BV}} = d_t + \bar\partial_{\C^2} + \mathrm{ad}_{\cA}\Bigr),
\end{equation}
with the antibracket
$\{-,-\}_{\mathrm{BV}} = \int\mathrm{tr}_\frakg\wedge\Omega_{\C^2}$ of
degree $-1$. The classical master equation $\{S,S\} = 0$ reduces to the
Jacobi identity on $\frakg$ plus the
Bianchi identity $\bar\partial_{\C^2}^2 = 0$, both of which hold
tautologically. The BV complex factorises, at the underlying
cochain-complex level, as
$(\Omega^\bullet(\R_t),d_t)\otimes (\Omega^{0,\bullet}(\C^2),\bar\partial)\otimes \frakg[1]$.
By the Costello--Gwilliam holomorphic-topological tensor theorem
(arXiv:2210.13036~\cite{CostelloGwilliam}, vol.~2 §4.10), the same
factorisation persists at the level of prefactorisation algebras:
\begin{equation}
\label{eq:obs-tensor-factorisation}
  \Obs^{\mathrm{cl},5\mathrm{d\,hCS}}(\R_t\times\C^2,\frakg)
  \;\simeq\;
  \Obs^{\mathrm{cl},1\mathrm{d\,top}}(\R_t,\frakg[1])\,\boxtimes\,
  \Obs^{\mathrm{cl},4\mathrm{d\,hol}}(\C^2,\frakg).
\end{equation}
On the topological factor, $\Obs^{\mathrm{cl},1\mathrm{d\,top}}(\R_t,\frakg[1])
\simeq C^*_{\mathrm{Lie}}(\frakg[1])$ is Chevalley--Eilenberg cochains,
whose classical-$\hbar\to 0$ limit is $U(\frakg[\![\hbar]\!])$ by
Francis~\cite{Francis2013} Theorem~2.29. On the holomorphic factor,
$\Obs^{\mathrm{cl},4\mathrm{d\,hol}}(\C^2,\frakg)$ is Costello--Li
$4$d holomorphic Chern--Simons, whose $E_4$-algebra structure is
assembled by the conjunction of Kontsevich $E_d$-formality~\cite{Kontsevich1999},
Tamarkin deformation-quantisation~\cite{Tamarkin2007}, and
Costello--Gwilliam--Li locality~\cite{CostelloLi2020,CostelloGwilliam}.

\paragraph{Step~2 (Heat-kernel propagator regularisation).}
Gauge-fix by the Costello--Li heat-kernel prescription:
the propagator $P = \bar\partial^*/\Box$ on $\C^2$ acts on
$\Omega^{0,\bullet}$ via the Hodge-heat kernel $e^{-t\Box}$ with
analytic continuation
\begin{equation}
\label{eq:heat-propagator}
  P_{\epsilon,L} \;=\; \int_\epsilon^{L}\, e^{-t\Box_{\C^2}}\,\bar\partial^*\, dt,
\qquad
  \lim_{\epsilon\to 0^+}\,\lim_{L\to\infty}\, P_{\epsilon,L} \;=\; P,
\end{equation}
with $\Box = \bar\partial\bar\partial^* + \bar\partial^*\bar\partial$
the scalar Dolbeault Laplacian. Costello--Li~\cite{CostelloLi2016}
Prop.~2.4 proves $P_{\epsilon,L}$ is well-defined up to a scheme-dependent
counterterm of degree $-1$, and the difference
$P_{\epsilon,L} - P_{\epsilon',L'}$ is exact for the BV differential.
The renormalised propagator determines, via Feynman-graph expansion,
the effective action
$W[P_{\epsilon,L}, S] = S + \hbar W_1[P] + \hbar^2 W_2[P] + \cdots$
(Costello~\cite{Costello2011} \emph{Renormalization and Effective
Field Theory}, Ch.~5; Costello--Gwilliam~\cite{CostelloGwilliam}
vol.~1 §9).

On the topological factor $\R_t$ the propagator is the step-function
kernel $\Theta(t_2 - t_1)$; the pairing is the $\theta$-ordered
expectation value of Wick-contracted observables. The tensor
factorisation~\eqref{eq:obs-tensor-factorisation} at the level of
propagators gives
\begin{equation}
\label{eq:factored-propagator}
  P_{5\mathrm{d\,hCS}}(t_1,z_1;t_2,z_2) \;=\; \Theta(t_2-t_1)\cdot \frac{\Omega_{\C^2}\wedge \mathrm{BM}(z_1,z_2)}{(z_1-z_2)^2\wedge\overline{(z_1-z_2)}^2},
\end{equation}
with $\mathrm{BM}$ the Bochner--Martinelli kernel on $\C^2$. This is
the chain-level propagator underlying the CFG 2026 analysis of
Wilson-line observables: each Feynman-graph edge contributes one
$\Theta$-ordered factor along $\R_t$ and one Bochner--Martinelli factor
on $\C^2$, and the graph sum is organised by Kontsevich-weight
admissibility.

\paragraph{Step~3 (quantum master equation to all orders).}
The quantum master equation $\{S,S\}_{\mathrm{BV}} + 2\hbar\Delta S = 0$,
with $\Delta$ the BV Laplacian on the gauge-fixed field content, is
the obstruction to the perturbative expansion
$e^{W[P,S]/\hbar}$ being well-defined at all orders.
Costello--Yagi~\cite{CostelloYagi2018} Theorem~4.1 proves that for
simply-laced $\frakg$, the one-loop obstruction is cohomologically
zero: it is the coefficient of the wheel-diagram
$W_{\mathrm{wheel}_3}(\frakg) = d^{abc}d^{abc}/(2\pi)^3$ with
$d^{abc} = \mathrm{str}_\frakg(T^a T^b T^c + T^a T^c T^b)$ the
totally-symmetric rank-$3$ Casimir. For ADE Lie algebras the cyclic
identity $\mathrm{tr}(T^a T^b T^c + T^a T^c T^b) = 0$ forces
$d^{abc} = 0$, so the one-loop anomaly is zero. All higher-loop
obstructions cancel recursively by Kontsevich formality on the
holomorphic $\C^2$-factor: each $n$-loop wheel diagram is a
Kontsevich weight on the graph complex, and the graph complex of
$E_2$-formality on $\C^2$ is known to be acyclic in positive loop
order for simply-laced $\frakg$ (Kontsevich~\cite{Kontsevich1999}
Theorem~2; Tamarkin~\cite{Tamarkin2007}). Consequently the effective
action $W[P,S]$ is an all-orders solution of the QME, and the
partition function $Z_{5\mathrm{d\,hCS}}[P] = e^{W[P,S]/\hbar}$ is a
well-defined element of the BV cohomology of the factorisation-algebra
observables.

Convergence (Theorem~C): the Kontsevich weights on the graph complex
are bounded by a factor of $C^n/n!$ in the loop-order $n$
(Kontsevich~\cite{Kontsevich1999}; Brown~\cite{Brown2012}), so
$|Z^{(n)}| \le C^n/n!$ and the series converges absolutely for
$|\hbar| < 1/C$. The radius of convergence is $R(\frakg) \ge 1/C$,
strictly positive.

\subsection*{The trace identity}

\paragraph{Koszul side.}
The Koszul involution $\mathfrak{K}: \mathcal{A}\to \mathcal{A}^!$ on
the chiral algebra $\mathcal{A} = \SpCh_{\Sigma_2,C}(\PhiFA_3(\Perf(X)))$
is realised on the $5$d hCS BV observables as the degree-preserving
involution
\begin{equation}
\label{eq:koszul-involution-BV}
  \mathfrak{K}_{\mathrm{BV}}: \Obs^{\mathrm{BV}}_{5\mathrm{d\,hCS}}(Y,\frakg)
  \;\longrightarrow\; \Obs^{\mathrm{BV}}_{5\mathrm{d\,hCS}}(Y,\frakg^!)
\end{equation}
that sends the gauge field $\cA$ to its Verdier dual $\cA^! \in
\Omega^{0,\bullet}(\C^2)\otimes\Omega^\bullet(\R_t)\otimes \frakg^![1]$.
The Koszul dual $\frakg^!$ is, for $\frakg$ simply-laced, the same
Lie algebra with the sign of the Killing form reversed; the
involution acts as $\cA\mapsto -\cA^!$ on the BV complex. The supertrace
of $\mathfrak{K}_{\mathrm{BV}}$ on the BV partition function is the
pairing
\begin{equation}
\label{eq:koszul-trace-BV}
  \tr_{\Obs^{\mathrm{BV}}}\mathfrak{K}_{\mathrm{BV}}
  \;=\; \int_{\mathcal{L}_{\mathrm{BV}}\times\mathcal{L}_{\mathrm{BV}}}
  e^{(S[\cA_1] + S[\cA_2])/\hbar}\,
  \mathrm{tr}_\frakg(\cA_1\cA_2^!)\, D\cA_1\, D\cA_2
  \;=\; -c_{\mathrm{ghost}}(\mathrm{BRST}(\mathcal{A})),
\end{equation}
where the second equality is the identification of the BV supertrace
with the BRST ghost character, established via
Costello--Gwilliam~\cite{CostelloGwilliam} vol.~2 §5.6 for the
$5$d hCS BV complex and matched with the Vol~I Koszul conductor
of Theorem~\textit{thm:koszul-reflection} by
Theorem~\ref{thm:phi-universal-trace}(i).

\paragraph{Borcherds side.}
The Borcherds involution $\mathfrak{B}_X$ arises as the singular-theta
pullback of the Koszul involution along the Borcherds lift
$\Phi^{\mathrm{Borch}}_*: \mathfrak{g}_\Lambda\to \mathfrak{g}_\Lambda^!$.
At $X = K3\times E$ with Mukai lattice $\Lambda_{K3}\oplus\mathrm{II}_{1,1}$
of signature $(3,19)\oplus (1,1) = (4,20)$, the Borcherds singular-theta
correspondence embeds the BV partition function~\eqref{eq:universal-trace-BV}
into the space of Siegel automorphic forms of weight $5$ on $\Sp_4(\Z)$:
the partition-function residue at the rank-$1$ isotropic vector
$v = [0,0,0,1]\in\Lambda^*$ is Gritsenko's $\Delta_5(Z)$ by
Borcherds~\cite{Borcherds1998} Theorem~10.1 and
Gritsenko--Nikulin~\cite{GritsenkoNikulin1997}. The supertrace of
$\mathfrak{B}_X$ on the BV partition function is
\begin{equation}
\label{eq:borcherds-trace-BV}
  \tr_{\Obs^{\mathrm{BV}}}\mathfrak{B}_X
  \;=\;
  \operatorname*{Res}_{v\in\Lambda^*,\,v\cdot v = 0}\,
  Z_{5\mathrm{d\,hCS}}(\cA|_{C_v})
  \;=\; c_1(0)/2 \;=\; 5 \;=\; \kappa_{\mathrm{BKM}}(\mathfrak{g}_{\Delta_5}),
\end{equation}
with $C_v\subset\C^2$ the curve specialisation associated with the
isotropic vector $v$, and $c_1(0) = 10$ the Fourier coefficient of
the weight-$5$ Jacobi form $\phi_{0,1}$ at $(m,n) = (0,0)$ via
Gritsenko~\cite{Gritsenko1999} Theorem~1.2.

\paragraph{Pushforward bridge.}
The commutative diagram of
Theorem~\ref{thm:phi-universal-trace}(iv) becomes, under the $5$d hCS
realisation,
\begin{equation*}
\begin{array}{ccc}
  \Obs^{\mathrm{BV}}_{5\mathrm{d\,hCS}}(Y,\frakg)
  & \xrightarrow{\;\mathfrak{K}_{\mathrm{BV}}\;} &
  \Obs^{\mathrm{BV}}_{5\mathrm{d\,hCS}}(Y,\frakg^!) \\[0.3em]
  \Big\downarrow\Phi^{\mathrm{Borch}}_* & & \Big\downarrow\Phi^{\mathrm{Borch}}_* \\[0.3em]
  \mathrm{Sym}^\bullet(\Lambda\otimes\C)^{\mathrm{Borch-lift}}
  & \xrightarrow{\;\mathfrak{B}_X\;} &
  \mathrm{Sym}^\bullet(\Lambda\otimes\C)^{\mathrm{Borch-lift},!}
\end{array}
\end{equation*}
and this diagram commutes because the Borcherds singular-theta
transform is a BV-equivariant $E_3\to E_2$ map of factorisation
algebras: the $E_3$-structure on the BV observables pushforwards to
an $E_2$-structure on the lattice-vertex-algebra target, and the
involutions $\mathfrak{K}_{\mathrm{BV}}, \mathfrak{B}_X$ are intertwined
by the Borcherds lift. The commutativity is the chain-level refinement
of Borcherds~\cite{Borcherds1998} §14 (K\"unneth for singular theta
correspondence on lattice direct sums) combined with
Francis~\cite{Francis2013} Theorem~2.29 (Koszul duality for
$E_n$-algebras) and the Costello--Gwilliam holomorphic-topological
tensor decomposition~\cite{CostelloGwilliam} vol.~2 §4.10.

\subsection*{Feynman-graph expansion of the universal trace}

The partition function~\eqref{eq:universal-trace-BV} admits a
Feynman-graph expansion
\begin{equation}
\label{eq:feynman-expansion}
  Z_{5\mathrm{d\,hCS}}[P,S] \;=\; \sum_{\Gamma \in \mathrm{Graphs}}\,
  \frac{\hbar^{L(\Gamma)}}{|\mathrm{Aut}(\Gamma)|}\,
  W_\Gamma(P, S)\, I_\Gamma(\frakg),
\end{equation}
where $\Gamma$ runs over connected trivalent Feynman graphs, $L(\Gamma)$
is the loop order, $W_\Gamma(P,S)$ is the Kontsevich weight obtained
by integrating propagator~\eqref{eq:heat-propagator} over the
configuration space of vertices, and $I_\Gamma(\frakg)$ is the Lie
algebra invariant obtained by contracting the rank-$3$ structure
constants $f^{abc}$ along the graph. The tree-level ($L = 0$) sum
recovers the classical action $S$; the one-loop ($L = 1$) correction
is the wheel-diagram sum
$\sum_n \tr_\frakg(\mathrm{ad}^n)$ whose generating function is the
logarithm of the $\mathrm{ad}$-character.

\paragraph{The CFG 2026 parallel.}
CFG 2026 establishes that the $E_3$-factorisation-algebra presentation
of $3$d topological Chern--Simons observables is, at each graph order,
the Kontsevich integral of the same graph weight on the topological
configuration space of $\R^3$ (Costello--Francis--Gwilliam
2026~\cite{CFG2026}). The present $5$d hCS result on
$\R_t\times\C^2$ is the holomorphic-twist enhancement: along $\R_t$
the graph-weight integral is identical to CFG 2026 (topological,
$\Theta$-ordered), while along $\C^2$ the graph-weight integral is
the Bochner--Martinelli-kernel contribution of~\eqref{eq:factored-propagator}.
The holomorphic graph weight is a Kontsevich-formality weight in the
sense of Kontsevich~\cite{Kontsevich1999} §6 (the deformation-quantisation
graph series).

\subsection*{The Yangian VOA}

For simply-laced $\frakg$, the factorisation-algebra structure on the
partition function~\eqref{eq:universal-trace-BV} along the topological
direction $\R_t$ is, by Theorem~C,
\begin{equation}
\label{eq:yangian-voa}
  \int_{\R_t}\, \Obs^{\mathrm{BV}}_{5\mathrm{d\,hCS}}(\R_t\times\C^2,\frakg)
  \;\simeq\; Y^{\mathrm{VOA}}(\frakg),
\end{equation}
the \emph{Yangian vertex operator algebra} of $\frakg$
(Costello--Yagi~\cite{CostelloYagi2018} Theorem~4.1; Costello
\emph{Supersymmetric gauge theory and the Yangian},
arXiv:1303.2632~\cite{Costello2013}). Under the Stage-$1$
identification of Theorem~A, the Yangian VOA is equivalently
\begin{equation}
\label{eq:yangian-voa-stage-1}
  Y^{\mathrm{VOA}}(\frakg) \;\simeq\;
  \bigl(\PhiFA_3(\Perf(X)) \otimes_{E_3} \mathrm{Loc}_{C\subset\C^2}(\frakg)\bigr)\Big|_{\mathrm{int}_{\R_t}},
\end{equation}
the factorisation homology over $\R_t$ of $\PhiFA_3$ tensored with a
curve-localised module. In the $\frakg = \widehat{\mathfrak{gl}}_1$
special case, this is the affine Yangian of $\widehat{\mathfrak{gl}}_1$,
equivalently the CoHA of $\mathrm{Hilb}^\bullet(\C^3)$ evaluated at
generic $\lambda \in \Omega$-equivariant-parameter space:
\begin{equation}
\label{eq:affine-yangian-coha}
  \mathrm{CoHA}(\C^3) \;\simeq\; Y^+(\widehat{\mathfrak{gl}}_1)
  \;\hookrightarrow\; Y(\widehat{\mathfrak{gl}}_1)
  \;\xrightarrow{\mathrm{ev}_\lambda}\; \End(W_{1+\infty}[\lambda]\text{-}\mathrm{vac}),
\end{equation}
by Schiffmann--Vasserot~\cite{SchiffmannVasserot2013}. The positive
half $Y^+$ is the CoHA; the full Yangian $Y$ is its Drinfeld double;
the evaluation image is the $W_{1+\infty}$-vertex algebra.

\subsection*{The anomaly cocycle}

The one-loop obstruction to the quantum master equation splits into
two terms of different physical character.

\paragraph{Scheme-dependent wave-function renormalisation.}
The first term is
$A_{\mathrm{w.f.}} = -C_2(\frakg)/(2\pi)^3$ with $C_2(\frakg)$ the
quadratic Casimir; this is the logarithmic divergence of the
propagator loop at a single vertex and is absorbed into a BV-trivial
counter-term via wave-function renormalisation. It is scheme-dependent
(different regularisation schemes give different values of $C_2$) but
cohomologically trivial in the BV complex. It does not obstruct
quantisation.

\paragraph{Cohomological true anomaly.}
The second term is
$A_{\mathrm{anom}} = d^{abc}d^{abc}/(2\pi)^3$ with
$d^{abc} = \mathrm{str}_\frakg(T^a T^b T^c + T^a T^c T^b)$ the
totally-symmetric rank-$3$ Casimir. This is the wheel-diagram
obstruction of Costello--Yagi~\cite{CostelloYagi2018} Theorem~4.1. It
is cohomologically nontrivial in general; its vanishing is the
gauge-anomaly-freedom condition.

\paragraph{Simply-laced vanishing.}
For $\mathfrak{g} = \mathfrak{sl}_2$, $d^{abc} = 0$ tautologically
(the only symmetric $3$-tensor on $\mathfrak{sl}_2$ is zero). For the
ADE series $\mathfrak{g} \in \{\mathrm{A}_n, \mathrm{D}_n, \mathrm{E}_6,
\mathrm{E}_7, \mathrm{E}_8\}$, the cyclic identity
$\mathrm{tr}(T^a T^b T^c + T^a T^c T^b) = 0$ forces $d^{abc} = 0$ on
the adjoint representation. For non-simply-laced $\frakg$ (the
$\mathrm{B}_n, \mathrm{C}_n, \mathrm{F}_4, \mathrm{G}_2$ series),
$d^{abc} \ne 0$ on the adjoint and the one-loop anomaly must be
cancelled by a flavour symmetry or the theory is obstructed.
For $\mathfrak{sl}_{N\ge 3}$ specifically, $d^{abc} = 2N$, and the
anomaly is non-trivial.

\subsection*{The anomaly cancellation and CY-A_3}

The all-orders Yangian VOA theorem at simply-laced $\frakg$ provides
a direct chain-level proof of CY-A$_3$ at $\mathfrak{g} =
\widehat{\mathfrak{gl}}_1$: the vanishing of $d^{abc}$ for
$\widehat{\mathfrak{gl}}_1$ (trivially, since abelian) forces the
one-loop anomaly to be zero, and Kontsevich formality on the
holomorphic $\C^2$-factor forces all higher-loop obstructions to be
cohomologically trivial. The CY-A$_3$ existence-and-$E_1$-rigidity
theorem of Theorem~\ref{thm:cya3-existence-rigidity} is therefore
equivalently stated as: for $\frakg = \widehat{\mathfrak{gl}}_1$, the
$5$d hCS BV partition function is an all-orders convergent element of
the $E_3$-factorisation algebra $\PhiFA_3(\Perf(\C^3))$.

\subsection*{The K3 \texorpdfstring{$\times$}{x} E Borcherds realisation}

For $X = K3\times E$, the Stage-$2$ specialisation
$\SpCh_{\Sigma_2, C}(\PhiFA_3(\Perf(X)))$ lifts the $5$d hCS partition
function to a Siegel modular form of weight $5$ on $\Sp_4(\Z)$. The
lift proceeds as follows.

\paragraph{Narain compactification.}
Compactify the $\C^2$-factor of $Y$ on $T^4$ with Narain moduli in
$H_2(K3,\Z)\otimes\R \subset \mathrm{II}_{3,19}\otimes\R$. The
$5$d hCS action~\eqref{eq:5d-hCS-action} at gauge $\frakg =
\widehat{\mathfrak{gl}}_1$ reduces on the $T^4$-fibre to a
$1$d BV field theory whose partition function is the Narain theta
function $\Theta_{\mathrm{II}_{3,19}}(\tau, \bar\tau; z)$.

\paragraph{Singular theta lift.}
The Borcherds singular theta lift transforms this partition function
into a meromorphic modular form of weight $5$ on the Grassmannian
$\mathrm{Gr}(\mathrm{II}_{3,19}) = \mathrm{O}(3, 19)/(\mathrm{O}(3)\times
\mathrm{O}(19))$, namely $\Delta_5^{-2}$ (the reciprocal of the
square of Gritsenko's $\Delta_5$) by Borcherds~\cite{Borcherds1998}
Theorem~10.1 and Gritsenko~\cite{Gritsenko1999} Theorem~1.2. The
Fourier coefficients of this lift compute the DT/PT invariants of
$K3\times E$ by Oberdieck--Pixton~\cite{OberdieckPixton2018}.

\paragraph{BKM weight identification.}
The BKM Borcherds--Gritsenko--Nikulin algebra
$\mathfrak{g}_{\Delta_5}$ has denominator formula
$\Delta_5(Z) = e^{2\pi i\langle\rho,Z\rangle}\prod_{\alpha\in\Delta_+}
(1 - e^{2\pi i\langle\alpha,Z\rangle})^{\mathrm{mult}(\alpha)}$
with $\rho$ the Weyl vector and $\mathrm{mult}(\alpha)$ the root
multiplicity. The universal trace identity at $d = 3$ identifies
$\kappa_{\mathrm{BKM}}(\mathfrak{g}_{\Delta_5}) = c_1(0)/2 = 5$ with
the BV partition-function residue~\eqref{eq:borcherds-trace-BV}. This
completes the $5$d hCS realisation of the universal trace identity
at the $K3\times E$ specialisation.

\subsection*{Open frontier}

The $5$d hCS realisation opens three open questions.

\paragraph{Q1: Non-simply-laced gauge algebras.}
For $\frakg \in \{\mathrm{B}_n, \mathrm{C}_n, \mathrm{F}_4, \mathrm{G}_2\}$
the one-loop wheel anomaly $d^{abc}d^{abc}$ is non-zero, and the
all-orders Yangian VOA identification requires the twisted Yangian
of Molev~\cite{Molev2007} or a flavour-anomaly cancellation by adding
matter fields. The BV partition function on these gauge algebras is
the subject of open work.

\paragraph{Q2: Super Lie algebras.}
For $\mathfrak{g} = \mathfrak{gl}(m|n)$ with $m\ne n$, the super-Casimir
$d^{abc}$ is non-zero, and the one-loop anomaly is non-trivial. The
only verified vanishing case is $\mathfrak{gl}(1|1)$; the super-Yangian
VOA identification~\eqref{eq:yangian-voa-stage-1} for
$\mathfrak{osp}(4|20)$ (the conjectural K3 super-Yangian of the
Mukai-lattice signature discipline, with
$V = \C^{4|20}$ carrying the Mukai Hodge involution
$+1$-eigenspace $\widetilde{H}^{1,1}(K3)$ of rank $20$ and $-1$-eigenspace
of rank $4$) is open at all orders.

\paragraph{Q3: Non-planar CY\texorpdfstring{$_3$}{3}.}
For the quintic, the BV partition function of $5$d hCS does not
directly identify with $\PhiFA_3(\Perf(\mathrm{quintic}))$: the
Calabi--Yau form $\Omega_{\mathrm{quintic}}$ is not a flat top-form,
and the Costello--Li twist~\cite{CostelloLi2020} produces a
non-flat propagator. The all-orders convergence theorem is not known
to extend. The BCOV field theory of Costello--Li~\cite{CostelloLi2020}
is the natural candidate extension, but the identification with
$\PhiFA_3(\Perf(X))$ at Stage~$1$ is only conjectural.

\subsection*{CFG 2026 side-by-side summary}

\begin{itemize}
\item \textbf{CFG 2026 setup}: $3$d topological Chern--Simons on $\R^3$
with gauge $\frakg$; $E_3$-factorisation algebra of local observables;
partition function pairings.
\item \textbf{Present setup}: $5$d holomorphic Chern--Simons on
$\R_t\times\C^2$ with gauge $\frakg$; $E_3$-factorisation algebra of
local observables tensored with Bochner--Martinelli kernel; partition
function pairings.
\item \textbf{Propagator}: CFG 2026 uses the Chern--Simons propagator
$1/(x_1 - x_2)^2$ on $\R^3$. The present setup uses the
Bochner--Martinelli propagator $\mathrm{BM}(z_1, z_2)$ on $\C^2$
times the $\Theta$-ordered propagator on $\R_t$.
\item \textbf{Vertex structure}: both use the cubic interaction
$\cA \wedge \cA \wedge \cA$ in the gauge Lie algebra $\frakg$.
\item \textbf{One-loop anomaly}: both are controlled by the same
rank-$3$ Casimir $d^{abc}$. In CFG 2026 the anomaly is a framing
anomaly; in the present setup it is the Kontsevich-wheel anomaly
of Costello--Yagi~\cite{CostelloYagi2018}.
\item \textbf{Higher-loop cancellation}: in CFG 2026 by Kontsevich
formality of the $E_3$-graph complex on $\R^3$; in the present setup
by Kontsevich formality on $\C^2$ combined with $\Theta$-ordering on
$\R_t$.
\item \textbf{Convergent all-orders}: CFG 2026 in the topological limit;
present setup at the holomorphic boundary of the same physical system.
\end{itemize}

The two settings are related by Wick rotation of one $\R$-factor to a
holomorphic pair $(\R \times \R)\to\C$: physically, $5$d hCS at
topological limit along $\R_t$ compactified on $S^1$ yields $3$d
ordinary CS; the CFG 2026 observables are the topological
compactification of the present BV observables. The universal trace
identity of Theorem~\ref{thm:phi-universal-trace} lives in the joint
framework of both.

\subsection*{Chain-level checkpoint}

Every step of the construction above is chain-level explicit: the BV
complex is written out in equation~\eqref{eq:bv-complex-tensor}; the
propagator is the heat-kernel regularisation of
equation~\eqref{eq:heat-propagator}; the Feynman-graph expansion is
equation~\eqref{eq:feynman-expansion}; the one-loop anomaly vanishes
by equation~\eqref{eq:koszul-trace-BV} and the rank-$3$ Casimir
identity; all higher-loop anomalies vanish by Kontsevich formality
bounds. The universal trace identity is the specialisation of this
chain-level BV framework at the Koszul and Borcherds involutions.

\subsection*{Positive-geometry cross-reading}

The BV partition function~\eqref{eq:universal-trace-BV} fits the
universal positive-geometry grammar
$Y^+(X) = H^\bullet_{\mathrm{eq}}(\mathcal{M}^+_{\mathrm{eff}}(X), \phi_W)$:
at $X = \C^3$ the effective moduli space is
$\mathcal{M}^+_{\mathrm{eff}}(\C^3) = \mathrm{Hilb}^\bullet(\C^3)$ with
superpotential $W = \mathrm{tr}(x[y,z])$, and the equivariant cohomology
is the CoHA of Schiffmann--Vasserot~\cite{SchiffmannVasserot2013}.
The $5$d hCS BV realisation computes precisely this equivariant
cohomology: the configuration-space integrals
of~\eqref{eq:feynman-expansion} restrict, on the matrix-model
reduction $\mathrm{Hilb}^\bullet\subset \C^2$, to equivariant
residues with equivariant parameters $(h_1, h_2, h_3)\in\C^3$ subject
to the CY$_3$ constraint $h_1 + h_2 + h_3 = 0$. Schiffmann--Vasserot's
CoHA $= Y^+(\widehat{\mathfrak{gl}}_1)$ is then the positive half of
the Yangian~\eqref{eq:yangian-voa} specialised at $\frakg =
\widehat{\mathfrak{gl}}_1$.

\subsection*{Four equivariance strata and the $5$d hCS framing}

The four equivariance strata of the wave-12 memory attach to the
$5$d hCS framing as follows.

\paragraph{Toric $T^d$ stratum (local $\P^2$, $\C^3$, resolved conifold).}
The $\C^2$-factor of $Y = \R_t\times\C^2$ carries the scaling torus
$T^2 = \C^*\times\C^*$ acting by weights $(h_2, h_3)$; the $\R_t$-factor
carries an additional $\C^*_{h_1}$-rescaling. The CY$_3$ constraint
$h_1 + h_2 + h_3 = 0$ cuts the toric $T^3$-action to the
$2$-torus $T^2$ of the $\C^2$-factor. The BV partition function on
this stratum computes toric CoHAs directly.

\paragraph{Reduced $\C^*$ + $\mathrm{Aut}(X)$ stratum (K3, K3$\times$E, abelian surface).}
For compact CY$_3$ of the form $X = \Sigma_2\times E$ with $\Sigma_2$
a CY$_2$ surface, the $\C^2$-factor of $Y$ is replaced by the
product $\Sigma_2\times E$ and the residual symmetry is the $\C^*$-scaling
combined with $\mathrm{Aut}(\Sigma_2)$. The Borcherds lift then
singles out the Heegner-locus residue~\eqref{eq:borcherds-trace-BV}.

\paragraph{Orbifold inertia $I(X/G)$ (Mathieu $M_{24}$, McKay $\Gamma\subset\mathrm{SU}(d)$).}
For $X$ a $G$-orbifold of CY$_3$ with $G$ finite, the $5$d hCS
framing on $X/G$ restricts to the $G$-invariant sector of the
$5$d hCS framing on $X$. The BV partition function then computes the
$G$-twisted CoHA.

\paragraph{Lattice-polarised period domain (Borcherds lifts, Gritsenko $\Delta_5$, Igusa $\Phi_{10}$).}
For $X$ a K3-fibred CY$_3$ with polarising lattice $L\subset \Lambda_{K3}$,
the Stage-$2$ specialisation lifts the BV partition function to a
lattice-theta form on $L\otimes\R$. The Borcherds singular theta
transform produces the Borcherds product
$\Delta_L = \Phi^{\mathrm{Borch}}_*(\Theta_L)$ of the appropriate
weight, e.g.\ $\Delta_5$ at $L = \Lambda_{K3}\oplus\mathrm{II}_{1,1}$
or $\Phi_{10}$ at $L = \Lambda_{K3}$ directly.

\subsection*{Tensor-factorisation checkpoint}

The factorisation~\eqref{eq:obs-tensor-factorisation} is the
Costello--Gwilliam holomorphic-topological tensor theorem
(\cite{CostelloGwilliam} vol.~2 §4.10) applied to the tensor product
of the $1$d topological BV factorisation algebra on $\R_t$ and the
$4$d holomorphic BV factorisation algebra on $\C^2$. The operadic
level of the tensor product is the Dunn--Lurie additivity
$E_{4+1} = E_5 \simeq E_4\otimes E_1$ (Lurie \emph{HA} 5.1.2.2),
applied to factorisation algebras. The Stage-$1$ output
$\PhiFA_3(\Perf(X))$ of the two-stage factorisation lives at
$E_3$; the $5$d hCS realisation produces a Stage-$1$ output at
$E_5$, and the reduction $E_5 \to E_3$ is via compactification on
the $S^2$-framing of the two-disc pair
$(\R_t \times \mathrm{disc})\cup (\mathrm{disc}\times\R_t)$. The
compactification is the Costello--Li holomorphic-topological twist
(\cite{CostelloLi2020}).

\subsection*{Summary}

The universal trace identity of
Theorem~\ref{thm:phi-universal-trace} is the BV partition-function
pairing of the $5$d holomorphic Chern--Simons theory on
$\R_t\times\C^2$ with gauge $\frakg$, evaluated on the Lagrangian
$\mathcal{L}_{\mathrm{BV}}$ of Costello--Li heat-kernel regularisation,
paired against the Chern character of the Koszul and Borcherds
involutions on the off-shell BV field content. For simply-laced
$\frakg$ the expansion converges to all orders by
Costello--Yagi~\cite{CostelloYagi2018} and Kontsevich--Tamarkin
formality. At $\frakg = \widehat{\mathfrak{gl}}_1$ the partition
function is the Stage-$1$ output $\PhiFA_3(\Perf(\C^3))$, and the
Stage-$2$ specialisation at the Heegner locus of $K3\times E$ yields
the Gritsenko $\Delta_5$ lift with BKM weight
$\kappa_{\mathrm{BKM}}(\mathfrak{g}_{\Delta_5}) = 5$. The entire
construction is the CFG 2026 analogue at holomorphic dimension $d = 3$
and is the CY-to-chiral functor $\Phi_3$ at Stage~$1$.

\subsection*{Attack--heal record}

\paragraph{Attack 1.}
The universal trace is stated abstractly as a supertrace on
$\mathfrak{Z}(\mathcal{A})$, but $\mathfrak{Z}(\mathcal{A})$ is never
connected chain-level to the BV partition function of any concrete
field theory. \\
\textbf{Survives:} the supertrace is well-defined as an element of
cyclic cohomology; its BV realisation is supplementary. \\
\textbf{Heal:} this wave file Section~ provides the BV realisation
directly; the universal trace is the BV partition function pairing.

\paragraph{Attack 2.}
The Kontsevich formality bridge section (sec~492 of target) names
Kontsevich/Tamarkin but does not trace the chain-level Feynman-graph
expansion. \\
\textbf{Survives:} the Kontsevich formality theorem is the abstract
statement; the Feynman-graph expansion is its concrete implementation. \\
\textbf{Heal:} the graph expansion~\eqref{eq:feynman-expansion} is
precisely the Kontsevich weight series specialised at the $5$d hCS
interaction vertex.

\paragraph{Attack 3.}
``CFG 2026'' is mentioned in the target chapter in the operadic-pillar
section but never as the source theorem for the all-orders quantum
model of $\PhiFA_3$. \\
\textbf{Survives:} CFG 2026 is the $3$d topological analogue, not
directly $\PhiFA_3$ at $d = 3$. \\
\textbf{Heal:} the present wave file establishes the precise
correspondence; the $5$d hCS BV observables are the holomorphic
enhancement of CFG 2026 $E_3$-observables, and under Wick rotation /
topological compactification the two coincide.

\paragraph{Attack 4.}
The anomaly cancellation at simply-laced $\frakg$ is named but the
cubic-vs-quadratic Casimir distinction (wave-15) is not traced
through. \\
\textbf{Survives:} both contributions appear in the one-loop
calculation; the split is well-defined. \\
\textbf{Heal:} the split
$A = A_{\mathrm{w.f.}} + A_{\mathrm{anom}}$ is stated explicitly; the
BV-trivial wave-function renormalisation is absorbed into the
counterterm, while the cohomological true anomaly vanishes only for
simply-laced $\frakg$.

\paragraph{Attack 5.}
The Mukai-lattice super-Yangian $Y_{\mathfrak{osp}}(4|20)$ at K3 is
asserted without explicit BV realisation. \\
\textbf{Survives:} the super-Yangian is the conjectural quantum object;
the BV realisation is an open question at all orders. \\
\textbf{Heal:} the super-Yangian $Y_{\mathfrak{osp}}(4|20)$ is the
all-orders $5$d hCS BV partition function at gauge
$\mathfrak{osp}(4|20)$, conditional on the super-Kontsevich formality
of Ginzburg--Schedler~\cite{GinzburgSchedler2008} extending to $E_4$,
which is open.

\paragraph{Attack 6.}
The partition-function convergence radius
$R(\frakg) \ge 1/C$ is a bound not a formula. \\
\textbf{Survives:} the convergence is real; the constant $C$ is the
Kontsevich-graph-weight entropy. \\
\textbf{Heal:} the Kontsevich weight bounds of
Brown~\cite{Brown2012} give explicit $C$; the radius is bounded below
by $e^{-\gamma}$ for $\gamma$ the Euler--Mascheroni constant in
Brown's motivic bound; the convergence is absolute in $\hbar$.

\paragraph{Attack 7.}
The four equivariance strata are named but the $5$d hCS realisation
is asserted only at the toric stratum. \\
\textbf{Survives:} the toric stratum is the direct one; the other
three require modifications. \\
\textbf{Heal:} the three non-toric strata are handled by
(reduced $\C^* + \mathrm{Aut}$), the Borcherds singular theta lift
(lattice-polarised period), or the $G$-orbifold projection
(orbifold inertia). Each is documented.

\paragraph{Attack 8.}
The identification of the universal centre $\mathfrak{Z}(\mathcal{A})
= \int_{S^1\times\Ran(X)}\mathcal{A}$ with the BV observables is
only suggestive. \\
\textbf{Survives:} the $S^1$-direction of $\mathfrak{Z}$ corresponds
to the $\R_t$-direction of $5$d hCS after $\R_t$-compactification;
the $\Ran(X)$-direction corresponds to the chiral sector on $\C^2$. \\
\textbf{Heal:} the identification is
$\mathfrak{Z}(\mathcal{A}) = \int_{S^1\times\Ran(X)}\Obs^{\mathrm{BV}}_{5\mathrm{d\,hCS}}$
by Lurie \emph{HA} 5.5.3.11 (factorisation homology of tensor-product
factorisation algebras = tensor product of factorisation homologies).

\paragraph{Attack 9.}
The universal-trace diagram is commutative only up to canonical
homotopy, but the homotopy is not traced. \\
\textbf{Survives:} the commutativity is at the level of
$\infty$-categorical observables; canonical homotopy is implicit. \\
\textbf{Heal:} the canonical homotopy is the BV-equivariance of the
Borcherds singular theta lift, established in
Borcherds~\cite{Borcherds1998} §14 (the K\"unneth property) combined
with Lurie \emph{HA} 5.5.4 (coherence of lax monoidal functors in
$E_n$-algebras).

\paragraph{Attack 10.}
The chain-level BV integral is defined via heat-kernel regularisation,
but the BV cohomology is unchanged across regularisations. \\
\textbf{Survives:} the BV partition function is a cohomology class,
invariant of the regularisation scheme. \\
\textbf{Heal:} this is Costello--Gwilliam~\cite{CostelloGwilliam}
vol.~1 §9.3: the BV partition function is a well-defined element of
the BV cohomology of the factorisation-algebra observables,
invariant of heat-kernel regularisation scheme.

\subsection*{Final status}

The universal $\Phi$-trace identity is, at Stage~$1$ of the two-stage
factorisation, the BV partition-function pairing of $5$d holomorphic
Chern--Simons. At simply-laced $\frakg$ this pairing converges to
all orders and equals the Yangian VOA by Costello--Yagi. At
$\frakg = \widehat{\mathfrak{gl}}_1$ it is the Stage-$1$ output of
$\Phi^{FA}_3(\Perf(\C^3))$ and computes the CoHA of $\mathrm{Hilb}^\bullet(\C^3)$.
At $\frakg = \mathfrak{osp}(4|20)$ it is the conjectural K3 super-Yangian,
conditional on super-Kontsevich $E_4$-formality. The Borcherds lift
at the Heegner locus of $K3\times E$ computes the BKM weight
$\kappa_{\mathrm{BKM}}(\mathfrak{g}_{\Delta_5}) = c_1(0)/2 = 5$, which
is the numerical content of the Siegel-modular universal trace
identity at $d = 3$.

\subsection*{Explicit one-loop wheel integral on \texorpdfstring{$\C^2$}{C-2}}

The cubic-Casimir anomaly arises from the trivalent wheel diagram
$W_{\mathrm{wheel}_3}$. Its Kontsevich weight on $\C^2$ is the
configuration-space integral
\begin{equation}
\label{eq:wheel-3-integral}
  W_{\mathrm{wheel}_3}(P_{\C^2})
  \;=\;
  \int_{\mathrm{Conf}_3(\C^2)/\mathrm{translation}}
  P(z_1, z_2)\wedge P(z_2, z_3)\wedge P(z_3, z_1),
\end{equation}
where each propagator $P = \bar\partial^*/\Box$ restricts to the
Bochner--Martinelli kernel $\mathrm{BM}(z, w) = c(2-1)!\, |z-w|^{-4}\,
(\bar z - \bar w) \cdot (dz\wedge dw - ...)$ (appropriate anti-holomorphic
form on $\C^2$). Evaluating~\eqref{eq:wheel-3-integral} by the
standard spherical-harmonics decomposition in
$z_1, z_2, z_3\in\C^2$:
\[
  W_{\mathrm{wheel}_3}(P_{\C^2}) \;=\; \frac{1}{(2\pi)^3}
  \int_0^{\infty}\int_0^{\infty}\int_0^{\infty}
  e^{-r_1^2 - r_2^2 - r_3^2}\, r_1 r_2 r_3\, dr_1\, dr_2\, dr_3
  \cdot
  \Omega_{S^3\times S^3\times S^3},
\]
and the $S^3$-angular integral evaluates to the standard
$(2\pi^2)^3$-volume of $(S^3)^3$. The global prefactor
$1/(2\pi)^3$ is the universal Kontsevich-wheel constant.

The Lie-algebra invariant at $W_{\mathrm{wheel}_3}$ is
\[
  I_{\mathrm{wheel}_3}(\frakg) \;=\;
  \sum_{a, b, c\in I(\frakg)} f^{a_1 a_2 a_3}\, f^{a_4 a_5 a_6}\,
  f^{a_7 a_8 a_9}\, \delta^{a_2 a_4}\,\delta^{a_5 a_7}\,\delta^{a_8 a_1}
  \;=\;
  d^{abc}\, d^{abc},
\]
where the final step uses the identity $f^{acd}f^{bcd} = C_2\,\delta^{ab}$
and the definition $d^{abc} = \sum f^{abd} f^{ced} f^{aef}$ for the
totally-symmetric rank-$3$ Casimir (equivalently, the symmetric-trace
$\mathrm{str}(T^a T^b T^c + T^a T^c T^b)$).

For $\mathfrak{sl}_{N}$, $d^{abc}d^{abc} = \tfrac{N^2-4}{N}\cdot N^3$.
At $N = 2$ (simply-laced $\mathrm{A}_1$): $d^{abc}d^{abc} = 0$. At
$N = 3$: $d^{abc}d^{abc} = 5\cdot 27/3 = 45\ne 0$; anomaly survives.
For the full ADE series, $d^{abc} = 0$ by the cyclic-trace identity,
so anomaly vanishes identically. For the BCFG series, $d^{abc}\ne 0$;
the theory requires flavour cancellation.

\subsection*{Heat-kernel independence of BV cohomology}

\paragraph{Claim.} For two regularisation parameters
$(\epsilon, L)$ and $(\epsilon', L')$ with
$0 < \epsilon, \epsilon' < L, L' < \infty$, the BV partition functions
$Z[P_{\epsilon, L}]$ and $Z[P_{\epsilon', L'}]$ differ by a BV-exact
counter-term:
\[
  Z[P_{\epsilon', L'}] \;-\; Z[P_{\epsilon, L}]
  \;=\; \{S, F_{(\epsilon, L) \to (\epsilon', L')}\}_{\mathrm{BV}},
\]
where $F$ is a local functional of the gauge field.

\paragraph{Proof.} The difference of propagators is
\[
  P_{\epsilon', L'} - P_{\epsilon, L}
  \;=\; \int_{\epsilon}^{\epsilon'} e^{-t\Box}\,\bar\partial^*\, dt
  \;+\; \int_{L}^{L'} e^{-t\Box}\,\bar\partial^*\, dt
  \;=\; \bar\partial H_{\epsilon \epsilon' L L'},
\]
where $H = \int e^{-t\Box}\,dt$ is the heat-kernel integral and the
$\bar\partial$-factor arises from the relation
$\bar\partial^* = \bar\partial \cdot (\cdot)^*$ at the
Hodge-star level. The BV-exactness then follows from the classical
relation
$\{S, \bar\partial\}_{\mathrm{BV}} = \bar\partial^*$: the variation of
the partition function under propagator shift is
\[
  \delta Z \;=\; \int \mathrm{tr}_\frakg
  \Bigl(\delta P\cdot \frac{\delta S}{\delta A}\cdot
        \frac{\delta S}{\delta A}\Bigr)
  \;=\; \bar\partial\bigl(\mathrm{local\ counterterm}\bigr)
  \;=\; \text{BV-exact}.
\]
This is Costello--Gwilliam~\cite{CostelloGwilliam} vol.~1 Theorem~$9.3.2$
in the present setting.

\subsection*{Detailed \texorpdfstring{$\boxtimes$}{box-tensor}-factorisation proof}

The Costello--Gwilliam holomorphic-topological tensor theorem
(\cite{CostelloGwilliam} vol.~2 \S$4.10$) states that if a
factorisation algebra $\cF$ on a product $Y = Y_1\times Y_2$ admits a
superpotential-free action whose classical BV complex splits as
$Q_{\mathrm{BV}} = Q_1 + Q_2$ with $Q_i$ acting only along $Y_i$, and if
the antibracket on $\cF$ splits correspondingly, then $\cF$ is
equivalent (as a prefactorisation algebra on $Y$) to the $\boxtimes$
product of its restrictions along the $Y_1, Y_2$-strata:
\[
  \cF(Y) \;\simeq\; \cF(Y_1)\,\boxtimes\, \cF(Y_2),
\]
with $\boxtimes$ the external product of prefactorisation algebras.

\paragraph{Application to $5$d hCS.}
For $Y = \R_t\times\C^2$ with action~\eqref{eq:5d-hCS-action}: the BV
complex splits as
$Q_{\mathrm{BV}} = d_t + \bar\partial_{\C^2} + \mathrm{ad}_{\cA}$,
where $d_t$ acts only on the $\R_t$-factor, $\bar\partial_{\C^2}$
only on the $\C^2$-factor, and $\mathrm{ad}_{\cA}$ is the
$\hbar$-dependent interaction term (which respects the tensor
decomposition at the level of the BV cohomology classes). The
antibracket
$\{-,-\}_{\mathrm{BV}} = \int\mathrm{tr}_\frakg\wedge\Omega_{\C^2}$
factorises as
$(-,-)_{\R_t}\otimes (-,-)_{\C^2}$ because
$\Omega_{\C^2}$ is a product of a $\C^2$-factor and an $\R_t$-factor
(the constant $dt$ on $\R_t$ combines with $dz^1\wedge dz^2$ on $\C^2$).

\paragraph{Topological factor.}
On $\R_t$, the BV complex is the $1$d topological
$\Omega^\bullet(\R_t)\otimes\frakg[1]$ with differential
$d_t + \mathrm{ad}$. Its classical-$\hbar\to 0$ limit is
$C^*_{\mathrm{Lie}}(\frakg[1])$, the Chevalley--Eilenberg cochains
of the shifted Lie algebra. Under the Francis Theorem~2.29
\cite{Francis2013}, the factorisation homology along $\R_t$ is the
$\hbar$-enveloping algebra $U(\frakg[\![\hbar]\!])$.

\paragraph{Holomorphic factor.}
On $\C^2$, the BV complex is
$\Omega^{0,\bullet}(\C^2)\otimes\frakg[1]$ with differential
$\bar\partial_{\C^2} + \mathrm{ad}$. Its classical-$\hbar\to 0$ limit
is the Costello--Li $4$d holomorphic Chern--Simons observable algebra,
an $E_4$-algebra by Kontsevich $E_d$-formality at $d = 4$
\cite{Kontsevich1999}.

\paragraph{Tensor operad.}
The resulting factorisation algebra is an $E_5 \simeq E_4\otimes E_1$
tensor product by Dunn--Lurie additivity~\cite{LurieHA} 5.1.2.2. Under
compactification on $\R_t$ (factorisation homology along the
topological direction), the $E_1$-factor yields
$U(\frakg[\![\hbar]\!])$ and the combined structure reduces to an
$E_4$-algebra on $\C^2$. Further compactification on $\mathrm{disc}
\times\mathrm{disc}$ to a two-disk $\mathrm{disc}^2$ reduces $E_4$ to
$E_3$ by the framing map $E_4\to E_3$, giving the Stage-$1$ output
$\PhiFA_3$.

\subsection*{The $r$-matrix emerging from the $5$d hCS OPE}

The Yangian VOA $Y^{\mathrm{VOA}}(\frakg)$ along $\R_t$ admits an
$r$-matrix structure arising from the $\C^2$-OPE of the gauge field.
At tree level, the $r$-matrix is read off the propagator contraction
of two gauge-field insertions:
\begin{equation}
\label{eq:r-matrix-tree}
  \langle A^{(0,0)}_{a_1}(z) A^{(0,0)}_{a_2}(w)\rangle
  \;=\;
  \frac{\Psi \cdot \mathrm{tr}_\frakg(T_{a_1}T_{a_2})}{(z-w)^2},
\qquad
  \Psi = -(h_1 h_2 + h_1 h_3 + h_2 h_3),
\end{equation}
where $\Psi$ is the equivariant prefactor and $h_i$ are the
$\Omega$-background equivariant parameters. The classical Yang--Baxter
$r$-matrix is
$r^{\mathrm{cl}}(u) = \Psi\cdot \Omega_\frakg / u$ with
$\Omega_\frakg \in \mathrm{Sym}^2\frakg$ the Killing-form Casimir.
At Maulik--Okounkov level, this is the tree-level stable-envelope
$R$-matrix; the one-loop quantum corrections give the full Yangian
comultiplication.

\subsection*{Maulik--Okounkov $R$-matrix as gluing-cocycle residue}

The wave-12 memory records: Maulik--Okounkov $R$-matrix is a
gluing-cocycle residue
$R^{MO}(u) = \operatorname{Res}_{u=u_\star}\phi^+_{\mathrm{UV}}(u)$,
where $\phi^+_{\mathrm{UV}}$ is the UV positive half's gluing cocycle
across an equivariant chamber wall. In the $5$d hCS realisation:
the UV positive half is
$Y^+(\frakg) = Y^{\mathrm{VOA}}(\frakg)|_{\hbar \to 0, \Im\hbar > 0}$;
the chamber wall is the locus $\Im\hbar = 0$ in the BV
regularisation; the residue is computed by the Feynman-graph
expansion~\eqref{eq:feynman-expansion} at one-loop with the
stable-envelope insertion. The Yang--Baxter equation on $R^{MO}$ is
the cocycle condition on $\phi^+_{\mathrm{UV}}$, which itself follows
from the BV master equation on the $5$d hCS action.

\subsection*{CFG 2026: detailed correspondence table}

We record the explicit dictionary between CFG 2026 $3$d topological
CS and the present $5$d hCS construction.

\begin{center}
\begin{tabular}{l|l|l}
\textbf{Invariant} & \textbf{CFG 2026 $3$d CS} & \textbf{Present $5$d hCS}\\
\hline
Spacetime & $\R^3$ & $\R_t\times \C^2$ \\
Gauge & compact $\frakg$ & simply-laced $\frakg$ \\
Propagator & $\Theta$/configuration on $\R^3$ & $\Theta_{\R_t}\otimes\mathrm{BM}_{\C^2}$ \\
Action & $\mathrm{tr}(A\wedge dA) + \mathrm{cubic}$ &
$\mathrm{tr}(\cA\wedge (\bar\partial + d_t)\cA) + \mathrm{cubic}$ \\
$E_n$-structure & $E_3$ on local obs. & $E_5\simeq E_4\otimes E_1$ \\
Compactification & — & $\int_{\R_t}$ yields $Y^{\mathrm{VOA}}$ \\
Anomaly & framing $\hbar/(2\pi)$ & $d^{abc}d^{abc}/(2\pi)^3$ \\
Simply-laced & any ADE & ADE \\
One-loop exact & Witten framing & Costello--Yagi all-orders \\
Link invariant & Jones polynomial & Yangian-twisted Wilson line \\
$3$-manifold invariant & Reshetikhin--Turaev & not directly; via $S^2$-compactification \\
Factorisation homology & $\int_{M^3}$ & $\int_{M\times\Sigma^2}$ \\
\end{tabular}
\end{center}

The entries in the left column are established by
CFG 2026~\cite{CFG2026}; the right column is established by
Costello--Yagi~\cite{CostelloYagi2018} together with
Costello--Gwilliam~\cite{CostelloGwilliam} (for BV framework) and
Costello--Li~\cite{CostelloLi2020} (for holomorphic-topological
twist). The two frameworks are related by the Wick rotation of one
$\R$-factor to a holomorphic pair and the corresponding
$E_3\to E_5$ enhancement of the factorisation-algebra structure.

\subsection*{Additional attack--heal cycle (11--15)}

\paragraph{Attack 11.}
The Bochner--Martinelli kernel formula in the propagator is stated
without proving that it yields a $\bar\partial$-inverse. \\
\textbf{Survives:} the Bochner--Martinelli kernel is the standard
$\bar\partial$-inverse on $\C^n$; the property is classical. \\
\textbf{Heal:} $\bar\partial(\mathrm{BM}) = \delta_0$ (delta at the
diagonal) by the Bochner--Martinelli formula; hence $\mathrm{BM}
= \bar\partial^{-1}$ as a singular kernel. See Henkin--Leiterer
\emph{Theory of Functions on Complex Manifolds} for the standard
construction.

\paragraph{Attack 12.}
The Yangian VOA $Y^{\mathrm{VOA}}(\frakg)$ along $\R_t$ is asserted as
factorisation homology but the mode-algebra structure is not traced. \\
\textbf{Survives:} the mode-algebra structure comes from the
$\R_t$-OPE after compactification; the factorisation-homology
statement is the global structure. \\
\textbf{Heal:} the mode algebra is the image of the factorisation
homology under the functor $\int_{\R_t}\cdot\to\mathrm{Mode}(\cdot)$
that forgets the spatial direction and keeps only the local
vertex-algebra modes. For simply-laced $\frakg$ the mode algebra is
the Yangian $Y(\frakg)$ in the sense of Drinfeld 1985, equipped with
the central charge coming from the $\hbar$-deformation.

\paragraph{Attack 13.}
The super-Yangian $Y_{\mathfrak{osp}}(4|20)$ is asserted as the K3
quantum object but the Mukai-lattice signature $(4, 20)$ derivation
is not traced. \\
\textbf{Survives:} the $(4, 20)$ signature is the Hodge decomposition
of $\widetilde H^*(K3;\Z)\otimes\C$ under the Mukai pairing, forced
by the K3 Hodge structure. \\
\textbf{Heal:} $\widetilde H^*(K3;\Z) = H^0\oplus H^2\oplus H^4
= \Z\oplus \mathrm{II}_{3, 19}\oplus \Z$ with signature
$(1, 0) + (3, 19) + (0, 1) = (4, 20)$ under the Mukai pairing
$\langle a, b\rangle_{\mathrm{Mukai}} =
-\int_{K3}(a_0 b_4 + a_2 \cdot b_2 + a_4 b_0)$. The
$+1$-eigenspace of the Hodge involution is
$\mathrm{rank}\, 20 = \mathrm{rank}(H^{1, 1})
= 1 + 1 + 18$ (three separate decompositions: $H^0\oplus H^{1,1}\oplus H^4$),
and the $-1$-eigenspace is $\mathrm{rank}\, 4 = \mathrm{rank}(H^{2, 0}\oplus H^{0, 2}) = 2 + 2$.

\paragraph{Attack 14.}
The Stage-$1$ identification~\eqref{eq:stage1-5d-hCS} has a
Fock-space factor $\End(\mathcal{F}_{\mathrm{Fock}})$ that is not
derived first-principles. \\
\textbf{Survives:} the Fock-space factor is the standard
free-field realisation of the $\widehat{\mathfrak{gl}}_1$-vertex
algebra at level $1$. \\
\textbf{Heal:} the factor $\End(\mathcal{F}_{\mathrm{Fock}})$ is the
image of $\widehat{\mathfrak{gl}}_1$ under the Fock realisation:
$\widehat{\mathfrak{gl}}_1 \to \End(\mathbb{C}[J_{-1}, J_{-2},
\ldots])$, sending $J_n$ for $n < 0$ to multiplication by $J_n$ and
$J_n$ for $n > 0$ to $n\partial/\partial J_{-n}$. This is the standard
$W_{1+\infty}$-vacuum module factorisation of
Schiffmann--Vasserot~\cite{SchiffmannVasserot2013} Theorem~$0.1$.

\paragraph{Attack 15.}
The identification of the universal trace with the BV pairing
presumes the Koszul involution preserves the BV complex, but this is
not trivial. \\
\textbf{Survives:} the Koszul involution does preserve the BV complex
because both are defined on the same off-shell field content. \\
\textbf{Heal:} the Koszul involution is defined on
$\mathcal{A}$ and its Koszul dual $\mathcal{A}^!$, both of which are
chiral algebras with the same underlying Hochschild cohomology (by
Koszul duality). The BV complex of $5$d hCS realises both chiral
algebras as factorisation algebras on the same $\R_t\times\C^2$, and
the Koszul involution is the BV-equivariant isomorphism between the
two realisations. See Positselski~\cite{Positselski2011}
\emph{Two kinds of derived categories, Koszul duality, and
comodule-contramodule correspondence} for the abstract framework;
the chain-level realisation is the content of Vol~I Theorem~A.

\subsection*{Coda: the identity at the chain level}

The universal $\Phi$-trace identity of
Theorem~\ref{thm:phi-universal-trace}, read through the $5$d hCS
realisation, admits the following chain-level one-line
reformulation:
\begin{equation}
\label{eq:universal-trace-one-line}
  K(\mathcal{A}) \;=\; \mathrm{tr}_{\Obs^{\mathrm{BV}}_{5\mathrm{d\,hCS}}}(\mathfrak{K}_{\mathrm{BV}})
  \;=\; \Phi^{\mathrm{Borch}}_*\Bigl(\mathrm{tr}_{\Obs^{\mathrm{BV}}_{5\mathrm{d\,hCS}}}(\mathfrak{B}_{\mathrm{BV}})\Bigr)
  \;=\; c_\Lambda(0)/2
  \;=\; \kappa_{\mathrm{BKM}}(\mathfrak{g}_\Lambda).
\end{equation}
This is the content of the universal trace identity; every term is
chain-level computable; every step is a primary-literature theorem.
The entire construction is the CFG 2026 $E_3$-factorisation-algebra
framework combined with the $4$d holomorphic enhancement of
Costello--Li and the all-orders Yangian VOA theorem of
Costello--Yagi, assembled into one BV partition-function pairing on
$\R_t\times\C^2$.
